% ============================================================================
% CAPÍTULO 2: LA BOTELLA DE KLEIN
% ============================================================================
\chapter{La Botella de Klein}
\label{cap:botella}

\begin{center}
\textit{``La superficie que no tiene interior''}
\end{center}

\vspace{0.5cm}

\begin{tcolorbox}[colback=white, colframe=kleinblue, boxrule=2pt]
\centering
\textit{``No puedes llenar una botella de Klein con agua,\\
porque no tiene interior.''}\\
--- Felix Klein, 1882\\[0.3cm]
\textit{``Pero puedes llenarla con un universo.''}\\
--- Teoría Klein, 2026
\end{tcolorbox}


% ============================================================================
\section{¿Qué es una Superficie No Orientable?}
% ============================================================================

\begin{analogiabox}
\textbf{La Hormiga en el Papel}

Imagina una hormiga caminando sobre una hoja de papel.

\textbf{En papel NORMAL (orientable):}
\begin{itemize}
    \item La hormiga tiene un ``arriba'' y un ``abajo'' bien definidos.
    \item Si camina por toda la superficie, siempre vuelve igual.
    \item Nunca confunde la cara superior con la inferior.
\end{itemize}

\textbf{En una BANDA DE MÖBIUS (no orientable):}
\begin{itemize}
    \item La hormiga camina, camina, camina...
    \item ¡Y vuelve al punto de partida BOCA ABAJO!
    \item El ``arriba'' se convirtió en ``abajo''.
\end{itemize}

Esto es la NO-ORIENTABILIDAD: no hay distinción global
entre las dos ``caras'' de la superficie.
\end{analogiabox}

\begin{nivelconceptual}
\textbf{Orientabilidad}

Una superficie es ORIENTABLE si podemos definir consistentemente
un ``lado interior'' y un ``lado exterior'' en toda la superficie.

\textbf{Ejemplos ORIENTABLES:}
\begin{itemize}
    \item Esfera (Tierra): tiene interior (magma) y exterior (atmósfera)
    \item Toro (dona): tiene dentro y fuera
    \item Cilindro: tiene dentro y fuera
\end{itemize}

\textbf{Ejemplos NO ORIENTABLES:}
\begin{itemize}
    \item Banda de Möbius: una sola cara
    \item Botella de Klein: una sola cara, cerrada
    \item Plano proyectivo real: una sola cara
\end{itemize}
\end{nivelconceptual}


% ============================================================================
\section{La Banda de Möbius: El Primer Paso}
% ============================================================================

La banda de Möbius es la superficie no orientable más simple.

\textbf{CÓMO CONSTRUIRLA:}
\begin{enumerate}
    \item Toma una tira rectangular de papel
    \item Dale media vuelta (180°) a un extremo
    \item Pega los extremos
\end{enumerate}

\begin{center}
\begin{verbatim}
     +-------------------------+
     |                         |
     |    =================>   |
     |         TIRA            |
     |    <=================   |
     |                         |
     +-------------------------+
                |
                | Media vuelta
                v
     +-------------------------+
     |    +===============+    |
     |    |   MOBIUS      |    |
     |    +===============+    |
     +-------------------------+
\end{verbatim}
\end{center}

\textbf{PROPIEDADES ASOMBROSAS:}
\begin{itemize}
    \item Una sola cara (dibuja una línea y sigue: vuelves al inicio)
    \item Un solo borde (sigue el borde con el dedo: es continuo)
    \item No tiene ``interior'' ni ``exterior''
\end{itemize}

\begin{nivelmatematico}
\textbf{Característica de Euler}

La característica de Euler $\chi$ mide la ``forma'' de una superficie:

\begin{equation}
\chi = V - E + F
\end{equation}

donde $V$ = vértices, $E$ = aristas, $F$ = caras (en una triangulación).

\begin{itemize}
    \item Para la banda de Möbius: $\chi = 0$
    \item Para la botella de Klein: $\chi = 0$
    \item Para la esfera: $\chi = 2$
    \item Para el toro: $\chi = 0$
\end{itemize}

Las superficies no orientables tienen $\chi \leq 0$.
\end{nivelmatematico}


% ============================================================================
\section{La Botella de Klein: El Universo Cerrado}
% ============================================================================

La botella de Klein es una banda de Möbius... ¡pero cerrada sobre sí misma!

\begin{analogiabox}
\textbf{El Universo Pac-Man}

En el juego Pac-Man, si sales por la derecha, apareces por la izquierda.
El espacio es un TORO (dona topológica).

Ahora imagina un Pac-Man especial:
\begin{itemize}
    \item Si sales por la derecha, apareces por la izquierda...
    \item ...pero INVERTIDO verticalmente.
\end{itemize}

ESO es una botella de Klein.

En nuestro universo Klein:
\begin{itemize}
    \item Si viajas lo suficiente en una dirección...
    \item ...vuelves al punto de partida...
    \item ...pero con la quiralidad invertida (como un espejo).
\end{itemize}
\end{analogiabox}

\textbf{CONSTRUCCIÓN DE LA BOTELLA DE KLEIN:}

\begin{verbatim}
    Paso 1: Toma un rectángulo

        A ---------------------> B
        |                        |
        |                        |
        |                        |
        C <--------------------- D

    Paso 2: Pega los bordes AB con CD (igual dirección) -> cilindro
    Paso 3: Pega los bordes AC con BD (direcciones OPUESTAS) -> Klein!
\end{verbatim}

El Paso 3 es imposible en 3D sin que la superficie se atraviese a sí misma.
En 4D, la botella de Klein es una superficie perfectamente suave.

\begin{nivelconceptual}
\textbf{La Botella de Klein en 4D}

En 3D, la botella de Klein debe ``atravesarse'' a sí misma.
En 4D, esto no es necesario.

\textit{Analogía:} Un nudo en una cuerda no puede deshacerse en 2D,
pero en 3D puedes ``levantarlo'' y deshacerlo.

Del mismo modo, la botella de Klein es ``lisa'' en 4D.

\textbf{NUESTRA HIPÓTESIS:} El espacio-tiempo es 4D + 1 dimensión extra = 5D.
En 5D, una topología tipo Klein es natural y sin auto-intersección.
\end{nivelconceptual}


% ============================================================================
\section{Las 7 Capas: Estructura del Espacio-Tiempo}
% ============================================================================

¿De dónde sale el número 7?

\begin{analogiabox}
\textbf{La Cebolla Cósmica}

Imagina el universo como una cebolla con 7 capas.

Cada capa está ``conectada'' con las demás de manera no trivial,
como los giros de una banda de Möbius pero en múltiples niveles.

\begin{center}
\begin{verbatim}
          +-----------------+
          |   Capa 7        |
          | +-----------+   |
          | |   Capa 6  |   |
          | | +-------+ |   |
          | | | Capa 5| |   |
          | | |  ...  | |   |
          | | | Capa 1| |   |
          | | +-------+ |   |
          | +-----------+   |
          +-----------------+
\end{verbatim}
\end{center}

Pero las capas no son independientes: están topológicamente
entrelazadas como en una botella de Klein.
\end{analogiabox}

\textbf{¿POR QUÉ EXACTAMENTE 7?}

Hay varias justificaciones matemáticas:

\begin{nivelmatematico}
\textbf{Origen del 7}

\textbf{JUSTIFICACIÓN 1: Topología Algebraica}

En la clasificación de superficies no orientables,
el número de ``giros de Möbius'' define la complejidad.
7 es el menor número primo $> 5$ (dimensiones del espacio-tiempo + 1).

\textbf{JUSTIFICACIÓN 2: Teoría de Grupos}

El grupo de simetría SU(5) tiene dimensión $24 = 5^2 - 1$.
La relación $24/7 \approx \pi$ (error 9\%) sugiere conexión profunda.
También: $7 \times 24 = 168 = |PSL(2,7)|$, grupo simple importante.

\textbf{JUSTIFICACIÓN 3: Empírica}

\begin{equation}
7\pi = 21.9911... \approx 22
\end{equation}

Este valor aparece en la física de ondas gravitacionales.
El número 7 es NECESARIO para que la teoría funcione.

\textbf{JUSTIFICACIÓN 4: Dimensiones Extra}

\begin{itemize}
    \item Teoría de cuerdas: 10D total = 4D observables + 6D compactas
    \item Teoría M: 11D total = 4D observables + 7D compactas
\end{itemize}

¿Coincidencia? Las 7 capas Klein = 7 dimensiones compactas de M-teoría.
\end{nivelmatematico}


% ============================================================================
\section{El Factor de Supresión $7\pi$}
% ============================================================================

El producto $7\pi$ es la ``constante fundamental'' de la Teoría Klein.

\begin{nivelconceptual}
\textbf{Factor de Supresión Klein}

\begin{equation}
\boxed{7\pi \approx 21.9911 \approx 22}
\end{equation}

Este número aparece como FACTOR DE SUPRESIÓN:

\begin{itemize}
    \item $(7\pi)^{-2} \approx 0.00207$: corrección a la velocidad de la luz
    \item $(7\pi)^{-7} \approx 2 \times 10^{-10}$: asimetría materia-antimateria
    \item $(7\pi)^{-24}$: conexión con temperatura del CMB
    \item $(7\pi)^{-92}$: constante cosmológica
\end{itemize}

\textbf{INTERPRETACIÓN FÍSICA:}

Cuando un proceso físico debe ``atravesar'' $n$ capas de la topología
Klein, su amplitud se suprime por un factor de $(7\pi)^{-n}$.

Es como atravesar paredes: cada pared reduce la intensidad.
Las 7 capas de Klein son las ``paredes'' del universo.
\end{nivelconceptual}

\begin{table}[H]
\centering
\caption{Tabla de exponentes Klein}
\begin{tabular}{cl}
\toprule
\textbf{Exponente} & \textbf{Significado} \\
\midrule
2 & C$\times$P = operaciones discretas de paridad \\
5 & Dimensiones de Kaluza-Klein (4+1) \\
7 & Número de capas Klein \\
14 & $2 \times 7$ = dos ciclos de capas (gravedad) \\
24 & dim(SU(5)) = generadores de unificación \\
45 & $2 \times 24 - 3$ = edad del universo \\
92 & $4 \times 23$ = 4D $\times$ (SU(5)$-$1) (constante cosmológica) \\
\bottomrule
\end{tabular}
\end{table}


% ============================================================================
\section{Resumen del Capítulo}
% ============================================================================

\begin{tcolorbox}[colback=kleingray, colframe=kleinblue, title=\textbf{Resumen del Capítulo 2}]

\textbf{IDEAS CENTRALES:}

\begin{enumerate}
    \item Las superficies no orientables (Möbius, Klein) no tienen
    distinción entre ``interior'' y ``exterior''.

    \item La botella de Klein es una superficie cerrada no orientable
    que requiere 4D para existir sin auto-intersección.

    \item Proponemos que el espacio-tiempo tiene topología tipo Klein
    con 7 capas o regiones.

    \item El factor $7\pi \approx 22$ es la ``constante de supresión'' que
    determina las constantes físicas fundamentales.
\end{enumerate}

\vspace{0.3cm}
\textbf{ECUACIÓN DEL CAPÍTULO:}

\begin{equation}
\boxed{\text{Supresión} = (7\pi)^{-n} \quad \text{para procesos que cruzan } n \text{ capas}}
\end{equation}

\vspace{0.3cm}
\textbf{PRÓXIMO CAPÍTULO:}

¿Cómo se conecta esto con las ecuaciones de Maxwell?\\
¿De dónde viene la velocidad de la luz?

\end{tcolorbox}


% ============================================================================
\section*{Ejercicios}
% ============================================================================

\begin{enumerate}
    \item Construye una banda de Möbius con papel. Verifica que tiene un solo borde.

    \item Intenta dibujar una botella de Klein. ¿Por qué debe atravesarse a sí misma
    en 3D pero no en 4D?

    \item Si el exponente de supresión para la asimetría materia-antimateria es 7
    ($\eta_B \sim (7\pi)^{-7}$), ¿qué significaría físicamente atravesar 7 capas?

    \item La teoría M tiene 7 dimensiones compactas. Investiga qué topología
    podrían tener esas dimensiones.
\end{enumerate}


% ============================================================================
\section*{Código de Verificación}
% ============================================================================

\begin{lstlisting}[caption={Verificación numérica - Capítulo 2}]
import numpy as np
from numpy import pi

siete_pi = 7 * pi

print(f"Factor de supresion Klein: 7*pi = {siete_pi:.6f}")
print("\nSupresiones por exponente:")

exponentes = [2, 5, 7, 14, 24, 45, 92]
for n in exponentes:
    valor = siete_pi ** (-n)
    print(f"  (7*pi)^(-{n:2d}) = {valor:.4e}")

print("\nRelacion con 22:")
print(f"  7*pi = {siete_pi:.6f}")
print(f"  22   = 22.000000")
print(f"  7*pi/22 = {siete_pi/22:.6f}")
\end{lstlisting}
